% ============================================================================
%  ABSTRACT Y KEYWORDS
%
%  A&A usa abstract ESTRUCTURADO: \abstract toma CINCO argumentos entre llaves,
%  y los cinco son obligatorios como tokens, aunque el 1.o y el 5.o puedan ir
%  vacíos ({}).
%
%      1. Context      opcional  — trasfondo del problema
%      2. Aims         OBLIGATORIO — qué se propone el artículo
%      3. Methods      OBLIGATORIO — cómo se investigó
%      4. Results      OBLIGATORIO — qué se obtuvo
%      5. Conclusions  opcional  — qué se concluye en general
%
%  Los encabezados ("Context.", "Aims.", ...) los pone la clase; no escribirlos.
%
%  Si prefieres el abstract tradicional de un solo párrafo, usa la forma de un
%  único argumento: \abstract{...} — pero la estructurada es la preferida.
%
%  NO dejar líneas en blanco ni comentarios ENTRE los cinco grupos: una línea
%  vacía abre un párrafo nuevo y rompe la lectura de los argumentos. Por eso
%  van pegados abajo, y toda la explicación vive en este bloque.
%
%  Toda abreviatura de instrumento, método u observatorio se expande la primera
%  vez aquí y otra vez en la introducción:
%      ...very long baseline interferometry (VLBI)...
% ============================================================================

\abstract
{}
{Objetivo del trabajo.}
{Métodos empleados.}
{Resultados obtenidos.}
{}

% ----------------------------------------------------------------------------
%  KEYWORDS
%  Separadas por " -- " (en-dash rodeado de espacios).
%  Tomarlas TEXTUALMENTE de la lista oficial de A&A; no inventar.
%  Lista: https://www.aanda.org/for-authors/latex-issues/information-files
%
%  Candidatas para este proyecto (verificar en la lista antes de fijarlas):
%    black hole physics
%    gravitation
%    relativistic processes
%    celestial mechanics
%    methods: numerical
%    Galaxy: center
% ----------------------------------------------------------------------------
\keywords{black hole physics --
          gravitation --
          methods: numerical}
